\chapter{Sviluppo di applicazioni MapReduce in Hadoop}
\label{cap:hadoop-mapreduce}

Questo capitolo illustra come implementare programmi MapReduce utilizzando il framework Hadoop. L'obiettivo principale di Hadoop è astrarre la complessità dei sistemi distribuiti: il progettista o sviluppatore deve concentrarsi unicamente sulla definizione logica delle funzioni di \textit{Map} e \textit{Reduce}, mentre il framework gestisce autonomamente l'esecuzione distribuita, la parallelizzazione, la fase di \textit{shuffle and sort}, lo scheduling dei task e la sincronizzazione.

Le API di MapReduce in Hadoop sono scritte in Java. Sebbene in framework più recenti come Spark sia diffuso l'utilizzo di Python, per Hadoop la scrittura nativa e più performante avviene in Java.

\section{Architettura di un'applicazione MapReduce e Terminologia}

Un programma Hadoop MapReduce è costituito da tre componenti logiche fondamentali, ciascuna implementata tramite una specifica classe Java:
\begin{itemize}
    \item \textbf{Driver}: la classe che configura l'applicazione e ne gestisce il flusso di esecuzione.
    \item \textbf{Mapper}: la classe responsabile della fase di \textit{map}.
    \item \textbf{Reducer}: la classe responsabile della fase di \textit{reduce}.
\end{itemize}

Per comprendere a fondo il funzionamento del framework, è cruciale chiarire la terminologia utilizzata, in particolare la distinzione tra classi e istanze (oggetti):
\begin{itemize}
    \item \textbf{Job}: rappresenta l'esecuzione (o \textit{run}) completa del programma MapReduce su un determinato dataset.
    \item \textbf{Task}: è l'esecuzione di un singolo Mapper (Map task) o di un singolo Reducer (Reduce task) su una porzione specifica di dati. In un Job vi sono solitamente decine o centinaia di task eseguiti in parallelo.
    \item \textbf{Driver, Mapper, Reducer (senza il suffisso "class")}: indicano le singole \textit{istanze} (oggetti) delle rispettive classi create a runtime dal sistema. Ad esempio, il sistema crea numerosi oggetti Mapper, ognuno per analizzare una porzione diversa di dati.
    \item \textbf{Input split}: è una porzione di dimensione fissa dei dati di input. Logicamente, il file di input viene diviso in "split". Generalmente, uno split corrisponde più o meno alla dimensione di un blocco di HDFS (es. 128 MB). La corrispondenza non è perfetta al 100\% perché lo split deve tenere conto dei record (es. righe di testo) che iniziano in un blocco e finiscono nel blocco successivo. 
\end{itemize}

\subsection{Il livello di parallelizzazione}
Il numero di istanze della classe Mapper (i \textit{Mapper task}) creati dal sistema è \textbf{esattamente uguale al numero di input split}. 
Se un file è diviso in 10 blocchi su HDFS, Hadoop creerà 10 istanze del Mapper, assegnando ciascuna a un server in cui risiede fisicamente quel blocco di dati per evitare trasferimenti di rete (concetto di \textit{data locality}). Questa rappresenta la massima parallelizzazione ottenibile senza spostare i dati.

\section{Componenti e Flusso dei Dati}

\subsection{Il Driver}
Il Driver contiene il metodo \texttt{main()}, che funge da punto di ingresso (entry point) dell'applicazione. Viene eseguito sulla macchina client (non sui nodi del cluster) e ha il compito di:
\begin{itemize}
    \item Leggere gli argomenti passati da riga di comando.
    \item Configurare il Job (definire i percorsi di input/output, i formati, e i tipi di dato previsti).
    \item Dichiarare le classi utilizzate per il Mapper e il Reducer.
    \item Sottomettere il Job al cluster Hadoop e attenderne il completamento.
\end{itemize}

\subsection{Il Mapper}
Il Mapper "implementa" la fase di map e viene eseguito direttamente sui nodi del cluster. All'interno della classe, il metodo fondamentale è \texttt{map(...)}. 
Esso riceve in input una coppia \((chiave, valore)\) letta dallo split e viene \textbf{invocato una volta per ogni singola coppia}. Analizza i dati e produce in output un insieme (da 0 a N) di coppie intermedie \((chiave, valore)\).
I risultati intermedi prodotti dal Mapper \textbf{non vengono salvati su HDFS}, ma nel file system locale (disco fisso) del server in cui il Mapper è in esecuzione, al fine di velocizzare le operazioni ed evitare la latenza di replicazione di HDFS.

\subsection{Il Reducer e il partizionamento dei dati (Shuffle and Sort)}
Tra la fase di map e la fase di reduce, il sistema raggruppa tutti i risultati intermedi. I dati vengono ordinati in base alla chiave e inviati attraverso la rete ai nodi che eseguono le istanze del Reducer.
I dati vengono distribuiti applicando una funzione di \textit{hash} sulla chiave: tutte le coppie che condividono la stessa chiave finiranno inevitabilmente alla stessa istanza del Reducer.

\begin{remark}[La scelta del numero di Reducer]
Mentre il numero di Mapper è scelto automaticamente dal sistema in base al numero di split, \textbf{il numero di Reducer deve essere esplicitamente specificato nel Driver dallo sviluppatore}.
Il sistema non può sapere a priori quante chiavi distinte produrrà la fase di map. Lo sviluppatore deve impostare un numero ragionevole per parallelizzare la fase di reduce.
\begin{itemize}
    \item \textit{Assegnare troppi Reducer}: Se si impostano 200 Reducer per un problema che genera solo 100 chiavi distinte (es. 100 parole uniche), il sistema creerà 200 istanze. Di queste, 100 rimarranno vuote e verranno distrutte subito, sprecando risorse e causando un sovraccarico (overhead) inutile per lo scheduling.
    \item \textit{La funzione di Hash e lo sbilanciamento (Data Skew)}: La funzione di hash standard cerca di distribuire le chiavi uniformemente tra i Reducer. Tuttavia, se una chiave ha una frequenza anomala (es. la parola "Paolo" compare milioni di volte, mentre "Luca" solo 4 volte), il Reducer a cui viene assegnato "Paolo" diventerà un collo di bottiglia e impiegherà molto più tempo degli altri. In questi rari casi, lo sviluppatore può definire una funzione di partizionamento personalizzata.
\end{itemize}
\end{remark}

Il metodo \texttt{reduce(...)} del Reducer viene \textbf{invocato una sola volta per ogni chiave distinta}, ricevendo come input la chiave e la lista iterabile di tutti i valori a essa associati. 
L'output finale prodotto dal Reducer viene, questa volta, scritto nel distributed file system (es. HDFS). Poiché HDFS replica ogni blocco (di default 3 copie) per tolleranza ai guasti, questa fase di scrittura finale genererà traffico di rete.

% FIGURA ORIGINALE DA RECUPERARE DALLE SLIDE
% Slide di provenienza: Da 16 a 22 (Animazione "MapReduce data flow with a single reducer") e 23-25 ("MapReduce data flow with multiple reducers")
% Tipo: diagramma strutturale
% Descrizione: Diagramma che mostra i nodi che leggono gli input split (blu), l'invio dei risultati intermedi ordinati attraverso la rete (frecce verso il/i reducer) e il consolidamento finale del Reducer.
% Punto di inserimento: In questa sezione che descrive l'interazione tra Mapper e Reducer.
\begin{figure}[H]
    \centering
    %\includegraphics[width=0.9\textwidth]{placeholder_data_flow.png}
    \caption{Flusso dei dati in MapReduce: generazione degli input split, elaborazione nei Mapper sui nodi locali, fase di shuffle and sort attraverso la rete, ed esecuzione dei Reducer.}
    \label{fig:mapreduce-data-flow}
\end{figure}

\section{Tipi di Dato (Data Types) in Hadoop}

I programmi Hadoop non utilizzano i tipi base nativi di Java (come \texttt{String}, \texttt{int}, \texttt{long}, \texttt{float}) come chiavi o valori direttamente nelle firme dei metodi \texttt{map} e \texttt{reduce}. Per questioni di efficienza nel trasferimento dei dati sulla rete (serializzazione/deserializzazione), Hadoop definisce delle classi \textit{wrapper} specifiche, dette \textbf{Writable}:
\begin{itemize}
    \item \texttt{org.apache.hadoop.io.Text}: corrisponde a \texttt{String}.
    \item \texttt{org.apache.hadoop.io.IntWritable}: corrisponde a \texttt{Integer}.
    \item \texttt{org.apache.hadoop.io.LongWritable}: corrisponde a \texttt{Long}.
    \item \texttt{org.apache.hadoop.io.FloatWritable}: corrisponde a \texttt{Float}.
\end{itemize}

Ogni classe utilizzata come \textbf{valore} deve implementare l'interfaccia \texttt{Writable}.
Ogni classe utilizzata come \textbf{chiave} deve implementare l'interfaccia \texttt{WritableComparable}. La chiave non solo deve essere serializzata, ma deve anche poter essere \textit{comparata} e \textit{ordinata} durante la fase cruciale di shuffle and sort, affinché le chiavi identiche vengano raggruppate.

\section{Formati di Input e Output}

Prima che i dati arrivino al Mapper, il file su HDFS deve essere diviso e tradotto in coppie chiave-valore. Questo è gestito dalla classe astratta \texttt{InputFormat}, che valida il formato del file, divide i dati in Input Split e fornisce un \texttt{RecordReader} per l'estrazione delle coppie.

Esistono implementazioni predefinite di \texttt{InputFormat} molto utilizzate, specialmente per i file testuali:

\begin{itemize}
    \item \textbf{\texttt{TextInputFormat}}: È il formato standard per i file di testo. Il file viene diviso per righe (usando il ritorno a capo). 
    \begin{itemize}
        \item \textbf{Chiave}: l'offset (posizione in byte) della riga all'interno del file. Il tipo di dato è \texttt{LongWritable}.
        \item \textbf{Valore}: il contenuto dell'intera riga. Il tipo di dato è \texttt{Text}.
    \end{itemize}
    Spesso, all'interno del metodo map, l'offset (la chiave) viene ignorato perché non contiene informazioni semanticamente utili per l'analisi.

    \item \textbf{\texttt{KeyValueTextInputFormat}}: Utilizzato per file in cui ogni riga contiene già implicitamente una coppia chiave-valore, divisa da un separatore (di default una tabulazione \texttt{\textbackslash t}).
    \begin{itemize}
        \item \textbf{Chiave}: il testo precedente al separatore. Il tipo di dato è \texttt{Text}.
        \item \textbf{Valore}: il testo successivo al separatore. Il tipo di dato è \texttt{Text}.
    \end{itemize}
\end{itemize}

\begin{remark}[Attenzione ai tipi di dato con KeyValueTextInputFormat]
Durante la lezione il professore ha sottolineato un errore comune: con \texttt{KeyValueTextInputFormat}, \textbf{sia la chiave che il valore sono di tipo \texttt{Text}}. Se il testo precedente al separatore è un numero (es. \texttt{10125\textbackslash tMister John}), la chiave generata sarà l'oggetto \texttt{Text} contenente la stringa "10125", non un intero. Bisogna prestare attenzione quando si scrive il codice del Mapper, per evitare errori di \textit{type mismatch}.
\end{remark}

\begin{itemize}
    \item \textbf{\texttt{SequenceFileInputFormat}}: Un formato per file binari sequenziali, ottimizzato per framework Big Data.
\end{itemize}

% FIGURA ORIGINALE DA RECUPERARE DALLE SLIDE
% Slide di provenienza: 40-41 ("Getting Data to the Mapper")
% Tipo: diagramma concettuale
% Descrizione: Mostra la trasformazione dal file su HDFS -> InputFormat -> InputSplits -> RecordReader -> Mapper (coppie K,V).
% Punto di inserimento: Alla fine della spiegazione degli InputFormat.
\begin{figure}[H]
    \centering
    %\includegraphics[width=0.85\textwidth]{placeholder_inputformat.png}
    \caption{Processo di acquisizione dei dati: l'InputFormat definisce gli split e il RecordReader genera le coppie chiave-valore passate al Mapper.}
    \label{fig:input-format}
\end{figure}

Parallelamente, la classe \texttt{OutputFormat} (es. \texttt{TextOutputFormat}) definisce come salvare i risultati su HDFS. Il formato di default crea file di testo in cui la chiave e il valore sono separati da un tab, rendendoli perfettamente leggibili da un successivo job configurato con \texttt{KeyValueTextInputFormat}.

\section{Struttura base del codice Hadoop: Word Count Example}

Analizziamo ora la struttura del codice Java attraverso il classico problema del conteggio delle parole (Word Count). L'obiettivo è calcolare le occorrenze di ogni parola in un file testuale.

\subsection{Implementazione del Driver}

Il Driver eredita da \texttt{Configured} e implementa \texttt{Tool}. Lo sviluppatore si concentra sul metodo \texttt{run()}, dove imposta il lavoro.

\begin{lstlisting}[language=Java, caption=Struttura del Driver per il Word Count]
public class WordCount extends Configured implements Tool {
    @Override
    public int run(String[] args) throws Exception {
        int numberOfReducers = Integer.parseInt(args[0]);
        Path inputPath = new Path(args[1]);
        Path outputDir = new Path(args[2]);

        Configuration conf = this.getConf();
        Job job = Job.getInstance(conf);
        job.setJobName("WordCounter");

        // Impostazione dei path
        FileInputFormat.addInputPath(job, inputPath);
        FileOutputFormat.setOutputPath(job, outputDir);

        // Classi per i formati
        job.setInputFormatClass(TextInputFormat.class);
        job.setOutputFormatClass(TextOutputFormat.class);

        // Collegamento delle classi e del Driver
        job.setJarByClass(WordCount.class);
        job.setMapperClass(WordCountMapper.class);
        job.setReducerClass(WordCountReducer.class);

        // Tipi di output del MAPPER (fondamentale specificarli)
        job.setMapOutputKeyClass(Text.class);
        job.setMapOutputValueClass(IntWritable.class);

        // Tipi di output del REDUCER
        job.setOutputKeyClass(Text.class);
        job.setOutputValueClass(IntWritable.class);

        // Numero di reducer desiderati
        job.setNumReduceTasks(numberOfReducers);

        return job.waitForCompletion(true) ? 0 : 1;
    }

    public static void main(String args[]) throws Exception {
        int res = ToolRunner.run(new Configuration(), new WordCount(), args);
        System.exit(res);
    }
}
\end{lstlisting}

Il metodo \texttt{setMapOutputKeyClass} e il corrispettivo per il value sono essenziali: il sistema non può ricavare automaticamente questi tipi generici a runtime (per via della \textit{type erasure} in Java). Non è invece necessario specificare esplicitamente i tipi in input al Reducer, poiché corrispondono esattamente all'output del Mapper.

\subsection{Implementazione del Mapper}

La classe estende \texttt{Mapper<ChiaveIn, ValoreIn, ChiaveOut, ValoreOut>}. Per l'input file testuale di default, le chiavi in ingresso sono \texttt{LongWritable} (offset) e i valori \texttt{Text} (riga di testo).

\begin{lstlisting}[language=Java, caption=Struttura del Mapper per il Word Count]
class WordCountMapper extends Mapper<
    LongWritable, Text,   // Input: Offset riga, Contenuto riga
    Text, IntWritable> {  // Output: Parola, Contatore (1)
    
    protected void map(LongWritable key, Text value, Context context) 
        throws IOException, InterruptedException {
        
        // Spezza la riga (value) in parole usando gli spazi come delimitatore
        String[] words = value.toString().split("\\s+");

        for(String word : words) {
            String cleanedWord = word.toLowerCase();
            // Emette (parola, 1) tramite l'oggetto context
            context.write(new Text(cleanedWord), new IntWritable(1));
        }
    }
}
\end{lstlisting}
L'oggetto \texttt{Context} rappresenta la connessione con l'esecuzione nel cluster e mette a disposizione il metodo \texttt{write} per emettere le coppie prodotte.

\subsection{Implementazione del Reducer}

La classe estende \texttt{Reducer<ChiaveIn, ValoreIn, ChiaveOut, ValoreOut>}. Nel nostro esempio, l'input riceve l'output del Mapper (\texttt{Text}, \texttt{IntWritable}).

\begin{lstlisting}[language=Java, caption=Struttura del Reducer per il Word Count]
class WordCountReducer extends Reducer<
    Text, IntWritable,   // Input: Parola, lista di conteggi (1)
    Text, IntWritable> { // Output: Parola, Somma totale

    protected void reduce(Text key, Iterable<IntWritable> values, Context context) 
        throws IOException, InterruptedException {
        
        int occurrences = 0;

        // Scorre la lista di valori per quella specifica chiave e li somma
        for (IntWritable value : values) {
            occurrences = occurrences + value.get();
        }

        // Emette la parola e la sua frequenza totale
        context.write(key, new IntWritable(occurrences));
    }
}
\end{lstlisting}

\section{Ottimizzazione delle Prestazioni: Il Combiner}
\label{sec:combiner}

Nel problema Word Count appena visto, se un Mapper trova la parola "Hadoop" mille volte nel suo blocco, produrrà mille coppie \texttt{(Hadoop, 1)} che verranno interamente inviate sulla rete verso il Reducer. Per limitare questa congestione, Hadoop permette di utilizzare un \textbf{Combiner} (o mini-reducer locale).

Il Combiner esegue una pre-aggregazione dei risultati intermedi all'interno dello stesso server del Mapper, prima della fase di \textit{shuffle}. Così, le mille coppie verranno aggregate localmente nella singola coppia \texttt{(Hadoop, 1000)} prima di viaggiare sulla rete.

\begin{remark}[Condizioni per l'uso del Combiner]
\textbf{La funzione applicata dal Combiner deve essere associativa e commutativa} (come la somma o la ricerca del massimo). La media matematica, ad esempio, non lo è. Se si applica una funzione non adeguata, il sistema non restituirà alcun errore, ma il risultato finale sarà semplicemente sbagliato.
Inoltre, \textbf{l'esecuzione del Combiner non è garantita}. Il framework decide a runtime se eseguirlo o meno (a seconda delle statistiche interne su disco/RAM). Pertanto, la logica dell'applicazione (e soprattutto del Reducer) deve essere costruita in modo che il risultato finale sia corretto a prescindere dall'esecuzione del Combiner. Questo implica anche che i tipi di dato in output dal Combiner debbano coincidere strettamente con i tipi di dato in output dal Mapper.
\end{remark}

\subsection{Perché nel Word Count si sommano i valori anziché contarli?}
Durante la lezione è stata posta l'attenzione sul perché nel Reducer (e nel Combiner) del Word Count si effettui la somma dei valori (\texttt{occurrences += value.get()}) e non si incrementi semplicemente un contatore (\texttt{occurrences++}) per ogni elemento dell'iteratore.
Senza il Combiner, ricevendo \texttt{(Hadoop, [1, 1, 1])}, iterare e contare con \texttt{+1} produrrebbe correttamente il risultato 3. Ma se il Combiner entra in azione, il Reducer potrebbe ricevere \texttt{(Hadoop, [2, 1])}. Se si contassero solo gli elementi nell'iterabile, il Reducer restituirebbe erroneamente \texttt{(Hadoop, 2)}. Sommandone i valori, invece, si ottiene correttamente 3, garantendo la coerenza del programma indipendentemente dal passaggio per il Combiner.

\subsection{Implementazione e Dichiarazione}
Per abilitare il Combiner, bisogna istruire il Driver aggiungendo:
\begin{lstlisting}[language=Java]
job.setCombinerClass(WordCountCombiner.class);
\end{lstlisting}

Il Combiner estende la classe base astratta \texttt{Reducer}. Poiché nella quasi totalità dei casi (come nel Word Count) l'operazione da compiere localmente è identica a quella globale (una somma parziale vs una somma totale), si usa comunemente \textbf{la stessa classe del Reducer per implementare il Combiner}. Si può quindi omettere la creazione della classe \texttt{WordCountCombiner} e indicare:
\begin{lstlisting}[language=Java]
job.setCombinerClass(WordCountReducer.class);
\end{lstlisting}
Questo riduce il codice da mantenere ed è la pratica del 99\% delle applicazioni Hadoop.

% FIGURA ORIGINALE DA RECUPERARE DALLE SLIDE
% Slide di provenienza: 80 e 81 ("Word count example without combiner" e "Word count example with combiner")
% Tipo: Diagrammi di flusso paralleli
% Descrizione: Due schemi che confrontano l'output K,V dal mapper al reducer, con e senza il blocco intermedio del Combiner che esegue un'aggregazione locale, riducendo l'arità dei dati trasferiti.
% Punto di inserimento: Fine della spiegazione sul Combiner.
\begin{figure}[H]
    \centering
    %\includegraphics[width=0.9\textwidth]{placeholder_combiner_comparison.png}
    \caption{Confronto del flusso dati: a sinistra l'invio diretto delle coppie, a destra l'effetto di pre-aggregazione introdotto dall'utilizzo di un Combiner a livello di nodo.}
    \label{fig:combiner-comparison}
\end{figure}

\section{Creazione di Tipi di Dato Personalizzati}

Spesso il valore da emettere è un dato complesso (es. due numeri, o un intero e una stringa). È necessario definire tipi di dato (\textit{Data Type}) personalizzati.
Se il tipo personalizzato funge da \textbf{valore}, deve implementare l'interfaccia \texttt{Writable}. Vanno ridefiniti obbligatoriamente i metodi \texttt{write(DataOutput out)} (per la serializzazione sulla rete) e \texttt{readFields(DataInput in)} (per la deserializzazione nel nodo di destinazione).

Esempio di classe complessa per immagazzinare una somma e un contatore (ad es. per il calcolo della media):
\begin{lstlisting}[language=Java]
public class SumAndCountWritable implements org.apache.hadoop.io.Writable {
    private float sum = 0;
    private int count = 0;
    
    // Metodi getter e setter omessi per brevita'

    @Override
    public void write(DataOutput out) throws IOException {
        out.writeFloat(sum);
        out.writeInt(count);
    }

    @Override
    public void readFields(DataInput in) throws IOException {
        sum = in.readFloat();
        count = in.readInt();
    }

    @Override
    public String toString() {
        return "sum=" + sum + ",count=" + count;
    }
}
\end{lstlisting}

\begin{remark}[Ordine di Lettura e Scrittura]
Un dettaglio fondamentale, evidenziato dal professore, riguarda l'implementazione dei metodi \texttt{write} e \texttt{readFields}. L'ordine in cui le variabili vengono serializzate in byte nel metodo \texttt{write} deve essere \textbf{scrupolosamente identico} all'ordine con cui vengono lette nel \texttt{readFields}.
La rete invia solo uno stream continuo di byte, senza metadati che descrivono i tipi interni. Se si scrive prima il \texttt{float} e poi l'\texttt{int}, e nel metodo di deserializzazione si invertissero le chiamate leggendo prima un \texttt{int} e poi il \texttt{float}, \textbf{Java non restituirebbe alcun errore o eccezione a runtime}. I byte verrebbero mal interpretati silenziosamente producendo un risultato totalmente incorretto.
\end{remark}

Il metodo \texttt{toString()} viene ridefinito perché l'\texttt{OutputFormat} di Hadoop (ad esempio il \texttt{TextOutputFormat}) invoca questo metodo al momento di scrivere il file finale. Senza di esso, per gli oggetti Java non primari, verrebbe stampato l'identificativo esadecimale del puntatore della cella di memoria, che risulta inutile ai fini dell'applicazione.

Se il tipo personalizzato deve essere usato come \textbf{chiave}, esso deve implementare l'interfaccia \texttt{WritableComparable}. Oltre ai metodi visti, si devono implementare:
\begin{itemize}
    \item \texttt{compareTo()}: essenziale per lo \textit{shuffle and sort}.
    \item \texttt{hashCode()}: fondamentale affinché la funzione di partizionamento assegni gli oggetti uguali allo stesso Reducer.
\end{itemize}

\section{Condivisione di Parametri e Contatori}

\subsection{Passare parametri al Mapper e al Reducer}
La firma dei metodi \texttt{map} e \texttt{reduce} è fissa. Per passare piccoli parametri configurabili (es. una soglia numerica fornita come argomento riga di comando), si sfrutta l'oggetto \texttt{Configuration}.
Nel \textbf{Driver}, il parametro viene memorizzato sotto forma di stringa \textbf{prima} di instanziare il Job:
\begin{lstlisting}[language=Java]
Configuration conf = this.getConf();
conf.set("soglia_temperatura", "50"); 
// ATTENZIONE: Questo va fatto prima di Job.getInstance(conf)
Job job = Job.getInstance(conf);
\end{lstlisting}

Nel \textbf{Mapper} (o Reducer), il parametro viene recuperato in sola lettura dall'oggetto contesto:
\begin{lstlisting}[language=Java]
String sogliaStr = context.getConfiguration().get("soglia_temperatura");
int soglia = Integer.parseInt(sogliaStr);
\end{lstlisting}
Tale approccio permette solo la condivisione di poche e piccole stringhe, non per scambiare grandi moli di dati.

\subsection{I Contatori (Counters)}
I \textbf{Counters} sono interi usati per calcolare statistiche globali (ad es. contare il numero di record malformati in cui manca un campo) senza dover scrivere ulteriori job MapReduce separati. Se il programma incontrasse una riga errata e lanciasse un'eccezione, l'intero Job fallirebbe sprecando tutto il tempo impiegato in precedenza. Identificare, isolare la riga ed incrementare un contatore di errore è un design pattern molto comune.

Si definiscono nel Driver tramite \texttt{enum} di Java (es. \texttt{COUNTERS.ERROR\_COUNT}), inizializzandosi implicitamente a zero.
Nei metodi \texttt{map} o \texttt{reduce}, si richiamano in questo modo:
\begin{lstlisting}[language=Java]
context.getCounter(COUNTERS.ERROR_COUNT).increment(1);
\end{lstlisting}

\begin{remark}[Natura incrementale dei Counter]
I metodi \texttt{map} e \texttt{reduce} in esecuzione sui nodi \textbf{possono solo invocare l'incremento} del Counter (\texttt{increment(N)}), non possono mai "impostarlo" (\texttt{set(10)}). Avendo istanze indipendenti del Mapper in esecuzione concorrente, non si può prevedere uno stato sincronizzato in tempo reale del valore del contatore. Hadoop raccoglierà gli incrementi da tutti i nodi asincronamente e calcolerà il valore globale.
Il valore finale definitivo del Counter si potrà quindi estrarre e leggere in sicurezza \textbf{solo nel Driver}, al termine dell'esecuzione del Job, sfruttando \texttt{job.getCounters().findCounter(...).getValue()}.
\end{remark}

\section{Job "Map-Only" e ottimizzazioni avanzate}

\subsection{Map-Only Job}
Alcuni algoritmi richiedono semplicemente il filtraggio del dataset o una mappatura 1:1, analizzando le righe in totale indipendenza. Un esempio è "restituisci il nome e l'età per gli utenti maggiori di 47 anni". 
In questi casi, non vi è necessità di aggregare informazioni definendo "gruppi", il che rende superflua la fase di Reduce e l'oneroso trasferimento sulla rete dello \textit{shuffle and sort}.
Per implementare un \textbf{Map-Only job}, occorre configurare esplicitamente il numero di reducer a zero all'interno del Driver:
\begin{lstlisting}[language=Java]
job.setNumReduceTasks(0);
\end{lstlisting}
L'output della fase di Map verrà direttamente memorizzato nel distributed file system (con file tipicamente nominati \texttt{part-m-00000} anziché \texttt{part-r-00000}).
Il professore ha inoltre fatto notare che azzerare il numero di reducer nel Driver è \textbf{un'eccellente tecnica di debugging}: permette allo sviluppatore di verificare facilmente l'output esatto emesso dal Mapper e isolare eventuali bug prima che passino alla logica del Reducer.

\subsection{In-Mapper Combiner: un passo oltre il Combiner}
Mentre il Combiner classico prevede l'istanziamento di un'ulteriore classe dedicata che comunica col Mapper, un'ottimizzazione più avanzata è l'\textbf{In-Mapper Combiner}. L'aggregazione viene svolta nativamente \textit{all'interno} del Mapper sfruttando variabili di istanza in memoria RAM e l'intero ciclo di vita della classe, attraverso l'overriding di due metodi aggiuntivi:
\begin{itemize}
    \item \texttt{setup(Context context)}: Viene invocato \textbf{una sola volta per ciascuna istanza del Mapper task} (all'inizio, analogo a un costruttore). Viene usato per inizializzare un dizionario associativo/array (\textit{hash map}) in memoria per mantenere i contatori locali di ciascuna chiave.
    \item \texttt{cleanup(Context context)}: Viene invocato \textbf{una sola volta al termine} dell'esecuzione del Mapper, per inviare (tramite \texttt{context.write()}) i risultati consolidati in rete.
\end{itemize}

Nel metodo \texttt{map()}, lo sviluppatore non emette più le coppie, bensì si limita ad aggiornare le strutture dati pre-inizializzate nel metodo \texttt{setup}.

\begin{remark}[Rischio dell'In-Mapper Combiner: Memoria limitata]
L'uso dell'In-Mapper Combiner annulla il costo di serializzazione e instanziamento logico tra Mapper e Combiner, risultando incredibilmente rapido. Tuttavia, tutto il raggruppamento avviene nella memoria Heap di Java. La memoria RAM standard di un'istanza Mapper ammonta spesso a 1GB.
Se il numero di "chiavi distinte" del blocco input elaborato è smisurato (e l'array associativo dovesse assumere dimensioni abnormi), si incorrerà in un inevitabile \textbf{Out Of Memory Error}, causando l'uccisione (kill) completa dell'applicazione da parte di Hadoop. L'In-Mapper Combiner andrebbe usato unicamente se c'è garanzia che la struttura statistica da mantenere in memoria occupi spazi gestibili dal framework.
\end{remark}





\section{Pattern di Ottimizzazione e Filtraggio in MapReduce}

Ora esploreremo alcune tecniche avanzate e design pattern molto comuni per l'ottimizzazione e la gestione dei job MapReduce. In particolare, analizzeremo come implementare un \textit{In-Mapper Combiner} per ottimizzare l'uso della rete e come strutturare un job di puro filtraggio (Map-Only) gestendo il passaggio dinamico di parametri dal Driver ai nodi di calcolo.

\section{Ottimizzazione tramite In-Mapper Combiner (Esercizio 9)}

Il problema del Word Count, che abbiamo già affrontato utilizzando una classe Combiner tradizionale, può essere risolto in modo ancora più efficiente utilizzando un approccio chiamato \textbf{In-Mapper Combiner}. 

L'obiettivo rimane lo stesso: dato in input un file testuale non strutturato, vogliamo calcolare il numero di occorrenze di ciascuna parola. Tuttavia, invece di delegare la pre-aggregazione a una classe Combiner separata, eseguiamo questa operazione direttamente all'interno della classe Mapper, risparmiando l'overhead legato all'istanziazione di ulteriori oggetti e riducendo il traffico di rete.

\subsection{Funzionamento dell'In-Mapper Combiner}
In un'implementazione MapReduce standard (o con un Combiner tradizionale), il metodo \texttt{map} emette una coppia chiave-valore \((word, 1)\) per ogni singola parola letta. Nel pattern In-Mapper Combiner, invece, modifichiamo il ciclo di vita del Mapper sfruttando i metodi di inizializzazione e chiusura:

\begin{enumerate}
    \item \textbf{Metodo \texttt{setup}}: Viene invocato una sola volta per ciascuna \textit{istanza} della classe Mapper (all'avvio della gestione di un input split). Qui inizializziamo una variabile locale, tipicamente un array associativo come una \texttt{HashMap}, che utilizzeremo per mantenere le statistiche locali (in questo caso, il conteggio parziale delle parole).
    \item \textbf{Metodo \texttt{map}}: Viene invocato per ogni riga di input. Invece di emettere coppie con il metodo \texttt{context.write}, estraiamo le parole e aggiorniamo le frequenze direttamente nella \texttt{HashMap} locale in memoria.
    \item \textbf{Metodo \texttt{cleanup}}: Viene invocato una sola volta al termine dell'elaborazione di tutte le righe assegnate a quella specifica istanza del Mapper. Qui iteriamo sulle voci della \texttt{HashMap} ed emettiamo le coppie \((word, local\_count)\).
\end{enumerate}

In questo modo, l'output generato alla fine della fase di Map è equivalente a quello che produrrebbe un Combiner tradizionale, ma ottenuto in modo più diretto. Il Reducer, di conseguenza, non subisce alcuna modifica e continuerà ad aggregare i conteggi parziali ricevuti per calcolare il risultato globale.

\subsection{Limiti e Gestione della Memoria Principale}
Come sottolineato dal professore, sebbene questo approccio sia più efficiente, richiede un'attenzione particolare all'uso della \textbf{memoria principale (RAM)}. 

Le variabili locali utilizzate per memorizzare le statistiche (la \texttt{HashMap}) devono essere sufficientemente piccole da poter risiedere nella memoria principale del server in cui è in esecuzione l'istanza del Mapper. Generalmente, a ogni istanza del Mapper è assegnato circa 1 GB di memoria principale.

È fondamentale chiedersi: \textit{qual è la dimensione massima di questa struttura dati nel caso peggiore?}
Nel problema del Word Count, la dimensione massima della \texttt{HashMap} corrisponde al \textbf{numero massimo di parole distinte} presenti nello split di input gestito dal Mapper.
Se analizziamo un documento in inglese, il vocabolario tipico è di circa 20.000 - 30.000 parole distinte. Memorizzare 30.000 stringhe associate a un numero intero richiede uno spazio irrisorio, ben al di sotto del gigabyte a disposizione; pertanto, in questo scenario non si verificheranno errori.

Tuttavia, se il dominio del problema prevedesse, ad esempio, 1 miliardo di parole o chiavi distinte all'interno di un singolo split, la \texttt{HashMap} saturerebbe la memoria, causando un \texttt{OutOfMemoryError} a runtime. È sufficiente che una sola istanza del Mapper vada in errore di memoria per causare il fallimento (kill) dell'intera applicazione MapReduce, costringendo a ricominciare l'elaborazione da zero.
Se si prevede che i dati non possano essere contenuti nella memoria principale, è obbligatorio scartare l'In-Mapper Combiner e utilizzare la soluzione standard basata sulla classe Combiner.

\subsection{Esempio di Dati (Toy Example)}
Consideriamo il seguente file di input, presentato nelle slide per l'esercizio 9:

\begin{verbatim}
Toy example
file for Hadoop.
Hadoop running
example.
\end{verbatim}

L'output finale, calcolato combinando i conteggi parziali, sarà la lista delle parole con le rispettive frequenze globali:
\begin{verbatim}
(example, 2)
(file, 1)
(for, 1)
(hadoop, 2)
(running, 1)
(toy, 1)
\end{verbatim}

\subsection{Analisi dell'Implementazione Java}

Analizziamo ora il codice Java che implementa l'In-Mapper Combiner per il Word Count.

\subsubsection{La Classe Driver}
Nella classe Driver, la configurazione è pressoché identica a quella di un normale Word Count, ma con una differenza fondamentale: \textbf{non viene specificata alcuna classe Combiner}.

\begin{lstlisting}[language=Java]
// Configurazione del job MapReduce
Job job = Job.getInstance(conf); 
job.setJobName("Exercise #9 - Word Count - in-mapper combiner");

// Set input/output format
job.setInputFormatClass(TextInputFormat.class);
job.setOutputFormatClass(TextOutputFormat.class);

// Set map class e i tipi di output del map
job.setMapperClass(MapperBigData.class);
job.setMapOutputKeyClass(Text.class);
job.setMapOutputValueClass(IntWritable.class);

// Set reduce class e i tipi di output del reduce
job.setReducerClass(ReducerBigData.class);
job.setOutputKeyClass(Text.class);
job.setOutputValueClass(IntWritable.class);

job.setNumReduceTasks(numberOfReducers);
\end{lstlisting}

\subsubsection{La Classe Mapper}
La classe Mapper contiene il cuore logico dell'In-Mapper Combiner.

\begin{lstlisting}[language=Java]
class MapperBigData extends Mapper<LongWritable, Text, Text, IntWritable> {

    // Variabile locale per mantenere le statistiche
    HashMap<String, Integer> wordsCounts;

    @Override
    protected void setup(Context context) { 
        // Inizializzazione della variabile all'avvio del Mapper
        wordsCounts = new HashMap<String, Integer>(); 
    }

    @Override
    protected void map(LongWritable key, Text value, Context context) 
            throws IOException, InterruptedException {
        Integer currentFreq;
        
        // Estrazione delle parole ignorando gli spazi multipli
        String[] words = value.toString().split("\\s+");

        for (String word : words) {
            String cleanedWord = word.toLowerCase(); // Normalizzazione
            currentFreq = wordsCounts.get(cleanedWord);

            if (currentFreq == null) { 
                // Prima occorrenza della parola in questo split
                wordsCounts.put(new String(cleanedWord), new Integer(1));
            } else { 
                // Incremento delle occorrenze
                currentFreq = currentFreq + 1;
                wordsCounts.put(new String(cleanedWord), new Integer(currentFreq));
            }
        }
        // NOTA: Nessuna invocazione di context.write() all'interno del map!
    }

    @Override
    protected void cleanup(Context context) 
            throws IOException, InterruptedException {
        // Al termine dell'elaborazione, emettiamo i dati pre-aggregati
        for (Entry<String, Integer> pair : wordsCounts.entrySet()) {
            context.write(new Text(pair.getKey()), new IntWritable(pair.getValue()));
        }
    }
}
\end{lstlisting}
Come si nota, i dati vengono emessi sulla rete verso i Reducer soltanto nel momento in cui il metodo \texttt{cleanup} viene chiamato dal framework Hadoop.

\subsubsection{La Classe Reducer}
Come anticipato, la classe Reducer non subisce alcuna alterazione rispetto al Word Count classico, poiché riceve in input coppie dello stesso tipo, pre-aggregate.
\begin{lstlisting}[language=Java]
class ReducerBigData extends Reducer<Text, IntWritable, Text, IntWritable> { 
    @Override 
    protected void reduce(Text key, Iterable<IntWritable> values, Context context) 
            throws IOException, InterruptedException {
        int occurrances = 0;
        for (IntWritable value : values) {
            occurrances = occurrances + value.get();
        }
        context.write(key, new IntWritable(occurrances));
    } 
}
\end{lstlisting}

\section{Filtraggio dei Dati e Passaggio Dinamico di Parametri (Esercizio 12)}

Un altro design pattern fondamentale riguarda il puro filtraggio dei dati basato su condizioni fornite dinamicamente dall'utente. 
Consideriamo un problema in cui riceviamo in input una collezione di file testuali strutturati contenenti i valori giornalieri di PM10 per un insieme di sensori. 

Il formato di ogni riga è:
\begin{verbatim}
sensorId,date\tPM10 value (ug/m3)\n
\end{verbatim}

L'obiettivo è restituire solo i record in cui il valore di PM10 è strettamente inferiore a una specifica \textbf{soglia fornita dall'utente come parametro a riga di comando}.

\subsection{Il Pattern di Filtraggio (Map-Only Job)}
Analizzando il problema logico, notiamo che:
\begin{itemize}
    \item Ogni riga può essere valutata indipendentemente dalle altre.
    \item Non c'è necessità di aggregare o raggruppare record provenienti da sensori o date differenti.
\end{itemize}

Quando l'elaborazione non richiede raggruppamento (grouping) né aggregazione, i Reducer sono inutili e rappresentano un overhead di rete e di computazione. La soluzione ottimale è utilizzare un \textbf{Map-Only Job}, impostando il numero di task reducer a zero (\texttt{job.setNumReduceTasks(0)}). 

In un Map-Only job, l'output emesso dal metodo \texttt{map} non passa attraverso la fase di shuffle and sort, ma viene scritto direttamente nell'HDFS.

Riguardo al formato dell'output, il professore ha illustrato diverse alternative implementative, tutte ugualmente valide ma con differenze formali:
\begin{enumerate}
    \item \textbf{Chiave e Valore estratti dal testo}: Possiamo emettere come chiave \texttt{sensorId,date} e come valore il \texttt{PM10}. Questa è la soluzione adottata nel codice finale dell'esercizio.
    \item \textbf{Chiave vuota e Valore linea intera}: Emettere una stringa vuota come chiave e l'intera riga originale come valore.
    \item \textbf{NullWritable e Valore linea intera}: È possibile utilizzare \texttt{NullWritable} come chiave. Si tratta di un tipo di dato predefinito in Hadoop che rappresenta il valore nullo in forma serializzata. È l'oggetto più piccolo definibile nel framework, persino più leggero di una stringa vuota, e risulta particolarmente utile quando la logica non richiede una vera e propria chiave da associare al valore emesso.
\end{enumerate}

\subsection{Passaggio di Parametri tramite l'Oggetto Configuration}
La soglia di filtraggio non è nota a priori (hardcoded), ma viene passata all'applicazione tramite gli argomenti (\texttt{args}). Questo pone un problema: la classe Driver riceve gli argomenti, ma è la classe Mapper (che viene eseguita su cluster remoti) a dover applicare il filtro.

Per condividere piccole quantità di informazioni dal Driver ai task (Mapper o Reducer), si utilizza l'oggetto \texttt{Configuration}. Il meccanismo prevede la definizione di una \textbf{proprietà personalizzata} (personalized property) all'interno del Driver, che i task potranno successivamente recuperare.

Le proprietà personalizzate supportano esclusivamente valori di tipo \texttt{String}, per cui anche un valore numerico come $21.5$ o $21$ dovrà essere scritto e letto come testo, e convertito (cast) solo dal Mapper.

\subsubsection{Attenzione: L'Ordine di Inizializzazione della Configuration}
Il professore ha richiamato l'attenzione su un errore logico estremamente comune e insidioso, dimostrato eseguendo l'applicazione passo-passo.

Se si definisce la proprietà personalizzata \textit{dopo} aver creato l'istanza del \texttt{Job}, in questo modo:
\begin{lstlisting}[language=Java]
Configuration conf = this.getConf();
Job job = Job.getInstance(conf); // Creazione del Job
conf.set("maxThreshold", args[2]); // ERRORE! Impostazione successiva
\end{lstlisting}

Il Mapper \textbf{non troverà} la proprietà. Questo accade perché l'invocazione di \texttt{Job.getInstance(conf)} crea il Job e contestualmente \textbf{genera una copia} dell'oggetto \texttt{Configuration}. Le modifiche successive fatte all'oggetto \texttt{conf} originale nel Driver non si rifletteranno nella copia passata ai Mapper.

In questo scenario di errore:
\begin{itemize}
    \item Il metodo \texttt{context.getConfiguration().get("maxThreshold")} nel Mapper restituirà \texttt{null}.
    \item La conversione \texttt{Float.parseFloat(null)} fallirà in modo silenzioso in background sui nodi.
    \item L'applicazione terminerà creando la cartella di output, ma questa \textbf{risulterà completamente vuota}, senza che sulla console vengano stampati errori o eccezioni visibili (se non si ispezionano i log approfonditi dei container).
\end{itemize}

La regola fondamentale è: \textbf{tutte le proprietà personalizzate devono essere definite sull'oggetto Configuration prima di istanziare l'oggetto Job}.

\subsection{Esempio di Dati}
Consideriamo l'input dell'esercizio 12 fornito sulle slide:
\begin{verbatim}
s1,2016-01-01   20.5
s2,2016-01-01   60.2
s1,2016-01-02   30.1
s2,2016-01-02   20.4
s1,2016-01-03   55.5
s2,2016-01-03   52.5
\end{verbatim}
Se passiamo al programma una soglia (\textit{threshold}) pari a $21$, l'output conterrà esclusivamente le righe relative ai valori $20.5$ e $20.4$, essendo gli unici strettamente minori di $21$:
\begin{verbatim}
s1,2016-01-01   20.5
s2,2016-01-02   20.4
\end{verbatim}

\subsection{Analisi dell'Implementazione Java}

Esaminiamo la soluzione implementativa per questo problema di filtraggio parametrico.

\subsubsection{La Classe Driver}
La classe Driver imposta il parametro e definisce il job come Map-Only. Sfrutta inolte il \texttt{KeyValueTextInputFormat}, un input format molto utile quando i dati in input sono separati da un tab (\texttt{\textbackslash t}). Questo formato divide in automatico la stringa prima del tab (la pone nella chiave) e la stringa dopo il tab (la pone nel valore).

\begin{lstlisting}[language=Java]
public int run(String[] args) throws Exception {
    Path inputPath = new Path(args[0]);
    Path outputDir = new Path(args[1]);

    Configuration conf = this.getConf();

    // 1. DEFINIZIONE DELLA PROPRIETA' (PRIMA DEL JOB)
    // Inviamo la soglia presa da riga di comando ai mapper
    conf.set("maxThreshold", args[2]);

    // 2. CREAZIONE DEL JOB (Crea una copia della configurazione)
    Job job = Job.getInstance(conf);
    job.setJobName("Exercise 12");

    FileInputFormat.addInputPath(job, inputPath);
    FileOutputFormat.setOutputPath(job, outputDir);

    job.setJarByClass(DriverBigData.class);

    // Utilizzo del KeyValueTextInputFormat
    job.setInputFormatClass(KeyValueTextInputFormat.class);
    job.setOutputFormatClass(TextOutputFormat.class);

    job.setMapperClass(MapperBigData.class);
    job.setMapOutputKeyClass(Text.class);
    job.setMapOutputValueClass(FloatWritable.class);

    // IMPOSTAZIONE MAP-ONLY JOB
    job.setNumReduceTasks(0);

    return job.waitForCompletion(true) ? 0 : 1;
}
\end{lstlisting}

\subsubsection{La Classe Mapper}
Anche in questo caso viene sfruttato il metodo \texttt{setup}. Piuttosto che richiedere e parsare l'oggetto \texttt{Configuration} per ogni singola riga analizzata dal metodo \texttt{map} (che introdurrebbe un'inefficienza notevole visto che la soglia è costante), il valore viene estratto una volta sola all'avvio del task.

\begin{lstlisting}[language=Java]
class MapperBigData extends Mapper<Text, Text, Text, FloatWritable> {

    float threshold;

    @Override
    protected void setup(Context context) { 
        // Recupero la soglia una volta sola per l'intero split
        // Viene letto come stringa e convertito in float
        threshold = Float.parseFloat(context.getConfiguration().get("maxThreshold")); 
    }

    @Override
    protected void map(Text key, Text value, Context context) 
            throws IOException, InterruptedException {
        
        float measure;
        
        // Il value contiene il valore PM10 grazie al KeyValueTextInputFormat
        measure = Float.parseFloat(value.toString());

        // Filtraggio: se rispetta la condizione, emettiamo il record
        if (measure < threshold) {
            context.write(new Text(key), new FloatWritable(measure));
        }
        // Se non rispetta la condizione, il metodo termina senza emettere nulla.
    }
}
\end{lstlisting}
Poiché i Reducer sono impostati a $0$, l'output emesso da \texttt{context.write} nel Mapper viene scritto direttamente in formato testuale all'interno dei file di output in HDFS (ad esempio \texttt{part-m-00000}), producendo esattamente il risultato atteso.